% Options for packages loaded elsewhere
\PassOptionsToPackage{unicode}{hyperref}
\PassOptionsToPackage{hyphens}{url}
\documentclass[
  ignorenonframetext,
]{beamer}
\newif\ifbibliography
\usepackage{pgfpages}
\setbeamertemplate{caption}[numbered]
\setbeamertemplate{caption label separator}{: }
\setbeamercolor{caption name}{fg=normal text.fg}
\beamertemplatenavigationsymbolsempty
% remove section numbering
\setbeamertemplate{part page}{
  \centering
  \begin{beamercolorbox}[sep=16pt,center]{part title}
    \usebeamerfont{part title}\insertpart\par
  \end{beamercolorbox}
}
\setbeamertemplate{section page}{
  \centering
  \begin{beamercolorbox}[sep=12pt,center]{section title}
    \usebeamerfont{section title}\insertsection\par
  \end{beamercolorbox}
}
\setbeamertemplate{subsection page}{
  \centering
  \begin{beamercolorbox}[sep=8pt,center]{subsection title}
    \usebeamerfont{subsection title}\insertsubsection\par
  \end{beamercolorbox}
}
% Prevent slide breaks in the middle of a paragraph
\widowpenalties 1 10000
\raggedbottom
\AtBeginPart{
  \frame{\partpage}
}
\AtBeginSection{
  \ifbibliography
  \else
    \frame{\sectionpage}
  \fi
}
\AtBeginSubsection{
  \frame{\subsectionpage}
}
\usepackage{iftex}
\ifPDFTeX
  \usepackage[T1]{fontenc}
  \usepackage[utf8]{inputenc}
  \usepackage{textcomp} % provide euro and other symbols
\else % if luatex or xetex
  \usepackage{unicode-math} % this also loads fontspec
  \defaultfontfeatures{Scale=MatchLowercase}
  \defaultfontfeatures[\rmfamily]{Ligatures=TeX,Scale=1}
\fi
\usepackage{lmodern}
\ifPDFTeX\else
  % xetex/luatex font selection
\fi
% Use upquote if available, for straight quotes in verbatim environments
\IfFileExists{upquote.sty}{\usepackage{upquote}}{}
\IfFileExists{microtype.sty}{% use microtype if available
  \usepackage[]{microtype}
  \UseMicrotypeSet[protrusion]{basicmath} % disable protrusion for tt fonts
}{}
\makeatletter
\@ifundefined{KOMAClassName}{% if non-KOMA class
  \IfFileExists{parskip.sty}{%
    \usepackage{parskip}
  }{% else
    \setlength{\parindent}{0pt}
    \setlength{\parskip}{6pt plus 2pt minus 1pt}}
}{% if KOMA class
  \KOMAoptions{parskip=half}}
\makeatother
\setlength{\emergencystretch}{3em} % prevent overfull lines
\providecommand{\tightlist}{%
  \setlength{\itemsep}{0pt}\setlength{\parskip}{0pt}}
\usepackage{bookmark}
\IfFileExists{xurl.sty}{\usepackage{xurl}}{} % add URL line breaks if available
\urlstyle{same}
\hypersetup{
  hidelinks,
  pdfcreator={LaTeX via pandoc}}

\author{\texorpdfstring{}{}}
\date{}

\begin{document}

\begin{frame}{Related Work}
\protect\phantomsection\label{sec:related_work}
This section situates our framework within five intersecting research
traditions: Active Inference foundations, meta-cognition in cognitive
science, predictive processing, AI safety and value alignment, and
cognitive security.

\begin{block}{Active Inference Foundations}
\protect\phantomsection\label{active-inference-foundations}
The Free Energy Principle, proposed by \cite{friston2010free}, provides
the variational foundation for Active Inference, positing that
biological systems maintain their organization by minimizing variational
free energy. Subsequent work formalized Active Inference on discrete
state-spaces \cite{dacosta2020active}, established the
epistemic-pragmatic decomposition of Expected Free Energy
\cite{friston2015active}, and demonstrated connections to homeostatic
regulation \cite{pezzulo2015active}. The comprehensive textbook by
\cite{parr2022active} consolidated these developments into a unified
treatment. Recent advances include scale-free formulations bridging
discrete and continuous domains \cite{friston2025scalefree}, EFE-based
planning as variational inference \cite{champion2025efeplanning}, and
integrations with projective simulation \cite{feps2025}. Our work
extends this tradition by examining the meta-level implications of
generative model specification---a dimension largely implicit in prior
treatments.
\end{block}

\begin{block}{Meta-Cognition in Cognitive Science}
\protect\phantomsection\label{meta-cognition-in-cognitive-science}
Meta-cognition---cognition about cognition---was first systematically
studied by \cite{flavell1979metacognition}, who distinguished
meta-cognitive knowledge from meta-cognitive regulation.
\cite{nelson1990metamemory} formalized the monitoring-control framework
distinguishing object-level processing from meta-level oversight, a
distinction that maps directly onto our Cognitive/Meta-Cognitive axis.
More recently, \cite{fleming2012metacognition} developed computational
models of meta-cognitive sensitivity, demonstrating that confidence
calibration is itself a measurable cognitive capacity. Metacognitive
architectures have been implemented in AI systems, including MetaMind
for social reasoning \cite{metamind2025} and the SofAI framework for
fast-slow-metacognitive thinking \cite{sofai2025}. Our quadrant
framework formalizes these distinctions within the Active Inference
formalism, providing mathematical specificity to the monitoring-control
hierarchy.
\end{block}

\begin{block}{Predictive Processing}
\protect\phantomsection\label{predictive-processing}
The predictive processing paradigm \cite{clark2013whatever} proposes
that brains are prediction machines that minimize prediction error
through hierarchical generative models. \cite{hohwy2013predictive}
developed this into a comprehensive theory of perception, action, and
cognition, emphasizing the role of precision-weighting in hierarchical
inference. Active Inference inherits and extends predictive processing
by adding action selection through EFE minimization
\cite{friston2015active}. Our Data/Meta-Data axis captures a dimension
that predictive processing theorists have discussed informally---the
distinction between processing sensory data and processing information
\emph{about} sensory reliability---and formalizes it within the
generative model framework.
\end{block}

\begin{block}{AI Safety and Value Alignment}
\protect\phantomsection\label{ai-safety-and-value-alignment}
The value alignment problem---ensuring AI systems pursue objectives
aligned with human values---has been articulated by
\cite{russell2019human} as a fundamental challenge for advanced AI.
\cite{amodei2016concrete} catalogued concrete safety problems including
reward hacking, scalable oversight, and distributional shift. Active
Inference offers a distinctive approach: rather than specifying reward
functions, the modeler specifies preference priors (\(C\) matrix) and
epistemic structures (\(A\), \(B\), \(D\) matrices), making value
specification transparent and inspectable \cite{sajid2022active}. The
meta-pragmatic perspective developed here extends this insight by
showing that framework specification provides principled value alignment
through explicit parameterization of what matters and what can be known.
\end{block}

\begin{block}{Cognitive Security}
\protect\phantomsection\label{cognitive-security}
Cognitive security addresses the protection of cognitive processes from
adversarial manipulation. \cite{wardle2017information} developed
influential frameworks for understanding information disorder including
misinformation, disinformation, and malinformation.
\cite{benkler2018network} analyzed how networked information ecosystems
create vulnerabilities in collective sense-making. In the Active
Inference context, \cite{constant2019regimes} examined how cultural
affordances and epistemic niches shape the generative models through
which agents interpret the world. Our contribution maps these concerns
onto the quadrant framework, showing that different types of cognitive
manipulation target different quadrants---from sensory data manipulation
(Q1) to framework-level subversion (Q4)---and that defense strategies
must operate at corresponding meta-levels.
\end{block}

\begin{block}{Contrast with Reinforcement Learning}
\protect\phantomsection\label{contrast-with-reinforcement-learning}
A key differentiator of our framework is the contrast with standard
reinforcement learning (RL). In RL, reward functions \(R(s,a)\) are
externally specified scalars, and the agent's task is to maximize
cumulative reward \cite{tschantz2020scaling}. Active Inference replaces
this with preference priors \(C\) embedded within a generative model,
enabling multi-dimensional value landscapes, epistemic drives,
and---crucially---the possibility of treating the specification itself
as a research variable. Our quadrant framework makes this meta-level
structure explicit and amenable to systematic analysis.
\end{block}
\end{frame}

\end{document}
